%%%%%%%%%%%%%%%%%
% This is an sample CV template created using altacv.cls
% (v1.7.2, 28 August 2024) written by LianTze Lim (liantze@gmail.com). Compiles with pdfLaTeX, XeLaTeX and LuaLaTeX.
%
%% It may be distributed and/or modified under the
%% conditions of the LaTeX Project Public License, either version 1.3
%% of this license or (at your option) any later version.
%% The latest version of this license is in
%%    http://www.latex-project.org/lppl.txt
%% and version 1.3 or later is part of all distributions of LaTeX
%% version 2003/12/01 or later.
%%%%%%%%%%%%%%%%

%% Use the "normalphoto" option if you want a normal photo instead of cropped to a circle
% \documentclass[10pt,a4paper,normalphoto]{altacv}

\documentclass[10pt,a4paper,ragged2e,withhyper]{altacv}
%% AltaCV uses the fontawesome5 and simpleicons packages.
%% See http://texdoc.net/pkg/fontawesome5 and http://texdoc.net/pkg/simpleicons for full list of symbols.

% Change the page layout if you need to
\geometry{left=1.25cm,right=1.25cm,top=1.5cm,bottom=1.5cm,columnsep=1.2cm}

% The paracol package lets you typeset columns of text in parallel
\usepackage{paracol}

% Change the font if you want to, depending on whether
% you're using pdflatex or xelatex/lualatex
% WHEN COMPILING WITH XELATEX PLEASE USE
% xelatex -shell-escape -output-driver="xdvipdfmx -z 0" sample.tex
\iftutex 
  % If using xelatex or lualatex:
  \setmainfont{Roboto Slab}
  \setsansfont{Lato}
  \renewcommand{\familydefault}{\sfdefault}
\else
  % If using pdflatex:
  \usepackage[rm]{roboto}
  \usepackage[defaultsans]{lato}
  % \usepackage{sourcesanspro}
  \renewcommand{\familydefault}{\sfdefault}
\fi

% Change the colours if you want to
\definecolor{SlateGrey}{HTML}{2E2E2E}
\definecolor{LightGrey}{HTML}{666666}
\definecolor{DarkPastelRed}{HTML}{450808}
\definecolor{PastelRed}{HTML}{8F0D0D}
\definecolor{GoldenEarth}{HTML}{E7D192}
\colorlet{name}{black}
\colorlet{tagline}{PastelRed}
\colorlet{heading}{DarkPastelRed}
\colorlet{headingrule}{GoldenEarth}
\colorlet{subheading}{PastelRed}
\colorlet{accent}{PastelRed}
\colorlet{emphasis}{SlateGrey}
\colorlet{body}{LightGrey}

% Change some fonts, if necessary
\renewcommand{\namefont}{\Huge\rmfamily\bfseries}
\renewcommand{\personalinfofont}{\footnotesize}
\renewcommand{\cvsectionfont}{\LARGE\rmfamily\bfseries}
\renewcommand{\cvsubsectionfont}{\large\bfseries}


% Change the bullets for itemize and rating marker
% for \cvskill if you want to
\renewcommand{\cvItemMarker}{{\small\textbullet}}
\renewcommand{\cvRatingMarker}{\faCircle}
% ...and the markers for the date/location for \cvevent
% \renewcommand{\cvDateMarker}{\faCalendar*[regular]}
% \renewcommand{\cvLocationMarker}{\faMapMarker*}


% If your CV/résumé is in a language other than English,
% then you probably want to change these so that when you
% copy-paste from the PDF or run pdftotext, the location
% and date marker icons for \cvevent will paste as correct
% translations. For example Spanish:
% \renewcommand{\locationname}{Ubicación}
% \renewcommand{\datename}{Fecha}


%% Use (and optionally edit if necessary) this .tex if you
%% want to use an author-year reference style like APA(6)
%% for your publication list
% % When using APA6 if you need more author names to be listed
% because you're e.g. the 12th author, add apamaxprtauth=12
\usepackage[backend=biber,style=apa6,sorting=ydnt]{biblatex}
\defbibheading{pubtype}{\cvsubsection{#1}}
\renewcommand{\bibsetup}{\vspace*{-\baselineskip}}
\AtEveryBibitem{%
  \makebox[\bibhang][l]{\itemmarker}%
  \iffieldundef{doi}{}{\clearfield{url}}%
}
\setlength{\bibitemsep}{0.25\baselineskip}
\setlength{\bibhang}{1.25em}


%% Use (and optionally edit if necessary) this .tex if you
%% want an originally numerical reference style like IEEE
%% for your publication list
\usepackage[backend=biber,style=ieee,sorting=ydnt,defernumbers=true]{biblatex}
%% For removing numbering entirely when using a numeric style
\setlength{\bibhang}{1.25em}
\DeclareFieldFormat{labelnumberwidth}{\makebox[\bibhang][l]{\itemmarker}}
\setlength{\biblabelsep}{0pt}
\defbibheading{pubtype}{\cvsubsection{#1}}
\renewcommand{\bibsetup}{\vspace*{-\baselineskip}}
\AtEveryBibitem{%
  \iffieldundef{doi}{}{\clearfield{url}}%
}


%% sample.bib contains your publications
\addbibresource{sample.bib}
% \usepackage{academicons}\let\faOrcid\aiOrcid
\begin{document}
\name{Pim Nelissen}
\tagline{Computational Scientist with a specialisation in Physics}
%% You can add multiple photos on the left or right
\photoR{2.3cm}{pim}
% \photoL{2.5cm}{Yacht_High,Suitcase_High}

\personalinfo{%
  % Not all of these are required!
  
  %\email{pim.nelissen@icloud.com}
  \email{pi0274ne-s@student.lu.se}
  \phone{+46 760 226 691}
  %\mailaddress{Lund,}
  \location{Lund, Sweden}
  %\homepage{www.homepage.com}
  % \twitter{@twitterhandle}
  % \xtwitter{@x-handle}
  \linkedin{pim-nelissen}
  \github{pim-n}
  % \orcid{0000-0000-0000-0000}
  %% You can add your own arbitrary detail with
  %% \printinfo{symbol}{detail}[optional hyperlink prefix]
  % \printinfo{\faPaw}{Hey ho!}[https://example.com/]

  %% Or you can declare your own field with
  %% \NewInfoFiled{fieldname}{symbol}[optional hyperlink prefix] and use it:
  % \NewInfoField{gitlab}{\faGitlab}[https://gitlab.com/]
  % \gitlab{your_id}
  %%
  %% For services and platforms like Mastodon where there isn't a
  %% straightforward relation between the user ID/nickname and the hyperlink,
  %% you can use \printinfo directly e.g.
  % \printinfo{\faMastodon}{@username@instace}[https://instance.url/@username]
  %% But if you absolutely want to create new dedicated info fields for
  %% such platforms, then use \NewInfoField* with a star:
  % \NewInfoField*{mastodon}{\faMastodon}
  %% then you can use \mastodon, with TWO arguments where the 2nd argument is
  %% the full hyperlink.
  % \mastodon{@username@instance}{https://instance.url/@username}
}

\makecvheader
%% Depending on your tastes, you may want to make fonts of itemize environments slightly smaller
% \AtBeginEnvironment{itemize}{\small}

%% Set the left/right column width ratio to 6:4.
\columnratio{0.57}
\begin{paracol}{2}
	
\cvsection{Education}

\cvevent{M.Sc.\ Computational Science, Physics}{Lund University}{Sep 2024 -- \textbf{\textit{ongoing}}}{}
{\LaTeXraggedright
	\cvtag{Numerical ODE/PDE solvers}
	\cvtag{MATLAB}
	\par}
\divider

\cvevent{B.Sc.\ Physics}{Lund University}{Sep 2020 -- Jun 2024}{}
\begin{itemize}
	\item My thesis involved machine learning applied to simulated MC data of radioactive decays from superheavy nuclei.
\end{itemize}
{\LaTeXraggedright
	\cvtag{Python}
	\cvtag{Monte Carlo}
	\cvtag{Deep Learning}
	\cvtag{C\#}
	\cvtag{SQL}
	\par}

% Start a 2-column paracol. Both the left and right columns will automatically
% break across pages if things get too long.

\medskip

\cvsection{Experience}

% \cvevent{Research assistant}{Lund University}{2025 -- \textbf{\textit{ongoing}}}{Lund, Sweden}
%\divider

\cvevent{Sales specialist}{Jumbo Supermarkten}{2018 -- 2020}{Oss, Netherlands}
\begin{itemize}
	\item Responsible for the entire product experience, with occasional leadership of a team of 3-5 people.
\end{itemize}

\cvevent{Sales Assistant}{Jumbo Supermarkten}{2015 -- 2018}{Oss, Netherlands}


\medskip

\cvsection{Volunteering}

\cvevent{System Administrator}{DFRI}{Jun -- Aug 2023}{Sweden}
\begin{itemize}
\item I did some Linux system maintenance and projects at a non-profit dedicated to digital rights and privacy, particularly working with Docker and Nextcloud.
\end{itemize}

\divider

\cvevent{Digital Literacy Instructor}{City Library Oss}{Sep 2013 --Jun 2014}{Oss, Netherlands}
\begin{itemize}
	\item I hosted an education program for promoting digital skills within the local community, particularly aimed at the elderly.
\end{itemize}

% use ONLY \newpage if you want to force a page break for
% ONLY the current column
\newpage

% \divider
% \printbibliography[heading=pubtype,title={\printinfo{\faFile*[regular]}{Journal Articles}},type=article]
% \divider
%\printbibliography[heading=pubtype,title={\printinfo{\faUsers}{Conference Proceedings}},type=inproceedings]

%% Switch to the right column. This will now automatically move to the second
%% page if the content is too long.
\switchcolumn

\cvsection{About me}

\begin{quote}
Enjoys learning in depth and in breadth. Avid reader, and likes to spend time outdoors. Strong numerical background in physics.
\end{quote}

\medskip

% \cvsection{Most Proud of}
%\cvachievement{\faTrophy}{Fantastic Achievement}{and some details about it}
%\divider
%\cvachievement{\faHeartbeat}{Another achievement}{more details about it of course}
%\divider
%\cvachievement{\faHeartbeat}{Another achievement}{more details about it of course}

% \cvsection{Strengths}

% Don't overuse these \cvtag boxes — they're just eye-candies and not essential. If something doesn't fit on a single line, it probably works better as part of an itemized list (probably inlined itemized list), or just as a comma-separated list of strengths.

% The `ragged2e` document class option might cause automatic linebreaks between \cvtag to fail.
% Either remove the ragged2e option; or 
% add \LaTeXraggedright in the paragraph for these \cvtag
% {\LaTeXraggedright
% \cvtag{Curious}
% \cvtag{Eye for detail}
% \cvtag{Avid reader}
% \par}

% \medskip

\cvsection{Workshops}

% \cvevent{Computational Nuclear Science and Engineering}{IAEA Online}{Feb-May 2025}{}

% \cvevent{Introduction to Bayesian Statistical Learning}{Jülich Supercomputing Centre}{Mar 2025}{}

% \cvevent{Bayesian Statistical Learning}{Jülich Supercomputing Centre}{24-28 Mar 2025}{}
\cvevent{Programming with Fortran}{Leibniz Supercomputing Centre}{11-13 Mar 2025}{}
\cvevent{Git, code testing and documentation}{EPFL CECAM, BioNT}{4-6 Feb 2025}{}

%\cvsection{Courses}

%{\LaTeXraggedright
%\cvtag{Scientific Python}
%\cvtag{Numerical methods for ODEs}
%\cvtag{Simulation Tools}
%\cvtag{Deep learning}
%\cvtag{Advanced Python}
%\cvtag{Climate Modelling}
%\cvtag{SQL}
%\cvtag{C\# Programming}
%\par}

\cvsection{Technical Skill Pie}

% Adapted from @Jake's answer from http://tex.stackexchange.com/a/82729/226
% \wheelchart{outer radius}{inner radius}{
	% comma-separated list of value/text width/color/detail}


\wheelchart{1.3cm}{0.3cm}{%
	10/3em/accent/\textbf{\underline{Python}},
	9/2em/accent/\textbf{\underline{Linux}},
	6/4em/accent!50/{MATLAB},
	5/2em/accent/\textbf{\underline{Git}},
	4/5em/accent!50/{TensorFlow},
	4/2em/accent/\textbf{\underline{SQL}},
	3/2em/accent!30/{C\#}
}

\medskip

\cvsection{Languages}

\cvskill{Dutch (Native)}{5}
\cvskill{English (IELTS Certified)}{4.5}
\cvskill{German}{2} %% Supports X.5 values.
\cvskill{Swedish}{1.5}

%% Yeah I didn't spend too much time making all the
%% spacing consistent... sorry. Use \smallskip, \medskip,
%% \bigskip, \vspace etc to make adjustments.

% \divider

\medskip

\cvsection{Publications}

%% Specify your last name(s) and first name(s) as given in the .bib to automatically bold your own name in the publications list.
%% One caveat: You need to write \bibnamedelima where there's a space in your name for this to work properly; or write \bibnamedelimi if you use initials in the .bib
%% You can specify multiple names, especially if you have changed your name or if you need to highlight multiple authors.
\mynames{P. Nelissen,}
%% MAKE SURE THERE IS NO SPACE AFTER THE FINAL NAME IN YOUR \mynames LIST

\nocite{*}

\printbibliography[heading=pubtype,title={\printinfo{\faBook}{Bachelor's Thesis}},type=book]

\newpage

% \cvsection{Referees}

% \cvref{name}{email}{mailing address}
% \cvref{Prof.\ Alpha Beta}{Institute}{a.beta@university.edu}
% {Address Line 1\\Address line 2}

\end{paracol}


\end{document}
